\begin{description}
\item[a)] Let us designate the (unknown) original limiting magnitude of this
telescope as $m_1$, and refer to the intensity (at the focal plane) as $I_1$.
The new limiting magnitude is $m_2$. Due to the additional intensity, which
is due to the better reflectivity of the mirror surfaces, the stellar
intensity which will result in the same limiting illumination at the focal
plane, will be proportionally smaller, with the ratio of the reflectivities
of the new and old mirror coatings:
\begin{equation*}
I_2 = \frac{\varepsilon_1\varepsilon_1}{\varepsilon_2\varepsilon_2}\,I_1
\tag{4.1}
\end{equation*}
\hfill(2 p)

The reflectivities should be multiplied twice, since the light coming from
the star, loses some energy twice, suffering reflection on each of the two
mirrors.

Now applying the well-known magnitude formula:
\begin{equation*}
m_1 - m_2 = -2.5\log\frac{I_1}{I_2}
          = -2.5\log\frac{\varepsilon_2^2}{\varepsilon_1^2}
\tag{4.2}
\end{equation*}
\hfill(2 p)

With numerical values:
\[ m_1 - m_2 = -0.16^{\mathrm{m}} \]
And thus:
\[ m_2 = m_1 + 0.16^{\mathrm{m}} \]
\hfill(1 p)

\item[b)] Yes, it is well appreciable by most human eyes. Therefore,
especially the deep-sky-object hunters are using higher reflectivity mirrors
in their telescopes. \hfill(2 p)

\item[c)] The similar way of thinking gives:
\begin{equation*}
m_2 = m_1 + 2.5\log\frac{I_1}{I_2}
    = m_1 + 2.5\log\frac{\varepsilon_2^2\varepsilon_3}{\varepsilon_1^3}
\tag{4.3}
\end{equation*}
\hfill(2 p)

With numerical values:
\[ m_1 - m_2 = -0.25^{\mathrm{m}} \]
And thus:
\[ m_2 = m_1 + 0.25^{\mathrm{m}} \]
\hfill(1 p)
\end{description}
